
\documentclass{article}
\usepackage{colortbl}
\usepackage{makecell}
\usepackage{multirow}
\usepackage{supertabular}

\begin{document}

\newcounter{utterance}

\centering \large Interaction Transcript for game `matchit\_ascii', experiment `different\_grid', episode 3 with Qwen3.6{-}35B{-}A3B{-}FP8{-}without{-}reasoning.
\vspace{24pt}

{ \footnotesize  \setcounter{utterance}{1}
\setlength{\tabcolsep}{0pt}
\begin{supertabular}{c@{$\;$}|p{.15\linewidth}@{}p{.15\linewidth}p{.15\linewidth}p{.15\linewidth}p{.15\linewidth}p{.15\linewidth}}
    \# & $\;$A & \multicolumn{4}{c}{Game Master} & $\;\:$B\\
    \hline

    \theutterance \stepcounter{utterance}  
    & & \multicolumn{4}{p{0.6\linewidth}}{
        \cellcolor[rgb]{0.9,0.9,0.9}{
            \makecell[{{p{\linewidth}}}]{
                \texttt{\tiny{[P1$\langle$GM]}}
                \texttt{Te osalete ühises arvamismängus.} \\
\texttt{Eesmärk on välja selgitada, kas see ruudustik ja teine pilt, mida ainult mina näen, on samad.} \\
\texttt{Palun kirjeldage kõigepealt oma ruudustikku. Seejärel annan mina oma kirjelduse ja me saame teineteisele esitada küsimusi ruudustike kohta, et välja selgitada, kas need on samad. Alustage nüüd oma lühikest ruudustiku kirjeldust märgiga "KIRJELDUS:", millele järgneb kirjeldus. Ärge lisage midagi muud. See on teie ruudustik:} \\
\texttt{▢ X X X ▢} \\
\texttt{X ▢ X ▢ X} \\
\texttt{X X ▢ X X} \\
\texttt{X X ▢ X X} \\
\texttt{X X ▢ X X} \\
            }
        }
    }
    & \\ \\

    \theutterance \stepcounter{utterance}  
    & \multicolumn{4}{p{0.6\linewidth}}{
        \cellcolor[rgb]{0.9,0.9,0.9}{
            \makecell[{{p{\linewidth}}}]{
                \texttt{\tiny{[P1$\rangle$GM]}}
                \texttt{KIRJELDUS: Ruudustik on 5x5. Esimeses reas on keskel X ja ääres tühjad ruudud. Teises reas on X esimeses, kolmandas ja viies veerus. Kolmandas, neljandas ja viiendas reas on X esimeses, teises, neljandas ja viiendas veerus, keskel on tühik.} \\
            }
        }
    }
    & & \\ \\

    \theutterance \stepcounter{utterance}  
    & & & \multicolumn{2}{p{0.3\linewidth}}{
        \cellcolor[rgb]{0.9,0.9,0.9}{
            \makecell[{{p{\linewidth}}}]{
                \texttt{\tiny{[GM$|$GM]}}
                \texttt{continue} \\
            }
        }
    }
    & & \\ \\

    \theutterance \stepcounter{utterance}  
    & & \multicolumn{4}{p{0.6\linewidth}}{
        \cellcolor[rgb]{0.9,0.9,0.9}{
            \makecell[{{p{\linewidth}}}]{
                \texttt{\tiny{[GM$\rangle$P2]}}
                \texttt{Te osalete ühises arvamismängus.} \\
\texttt{Eesmärk on välja selgitada, kas see ruudustik ja teine pilt, mida ainult mina näen, on samad.} \\
\texttt{Palun kirjeldage kõigepealt oma ruudustikku. Seejärel annan mina oma kirjelduse ja me saame teineteisele esitada küsimusi ruudustike kohta, et välja selgitada, kas need on samad. Alustage nüüd oma lühikest ruudustiku kirjeldust märgiga "KIRJELDUS:", millele järgneb kirjeldus. Ärge lisage midagi muud. See on teie ruudustik:} \\
\texttt{▢ ▢ ▢ ▢ X} \\
\texttt{▢ ▢ ▢ X X} \\
\texttt{▢ ▢ X X X} \\
\texttt{▢ X X X X} \\
\texttt{X X X X X} \\
            }
        }
    }
    & \\ \\

    \theutterance \stepcounter{utterance}  
    & & & \multicolumn{4}{p{0.6\linewidth}}{
        \cellcolor[rgb]{0.9,0.9,0.9}{
            \makecell[{{p{\linewidth}}}]{
                \texttt{\tiny{[GM$\langle$P2]}}
                \texttt{KIRJELDUS: Ruudustik on 5x5. X{-}tähed moodustavad kolmnurga, mis algab paremast ülemisest nurgast ja laieneb allapoole vasakule. Täpsemalt:} \\
\texttt{{-} 1. rida: 4 tühja ruutu, siis 1 X (positsioonil 1,5)} \\
\texttt{{-} 2. rida: 3 tühja ruutu, siis 2 X{-}d (positsioonidel 2,4 ja 2,5)} \\
\texttt{{-} 3. rida: 2 tühja ruutu, siis 3 X{-}d (positsioonidel 3,3, 3,4 ja 3,5)} \\
\texttt{{-} 4. rida: 1 tühi ruut, siis 4 X{-}d (positsioonidel 4,2 kuni 4,5)} \\
\texttt{{-} 5. rida: 5 X{-}d (kogu rida on X{-}tähtedega täidetud)} \\
            }
        }
    }
     \\ \\

    \theutterance \stepcounter{utterance}  
    & & & \multicolumn{2}{p{0.3\linewidth}}{
        \cellcolor[rgb]{0.9,0.9,0.9}{
            \makecell[{{p{\linewidth}}}]{
                \texttt{\tiny{[GM$|$GM]}}
                \texttt{continue} \\
            }
        }
    }
    & & \\ \\

    \theutterance \stepcounter{utterance}  
    & & \multicolumn{4}{p{0.6\linewidth}}{
        \cellcolor[rgb]{0.9,0.9,0.9}{
            \makecell[{{p{\linewidth}}}]{
                \texttt{\tiny{[GM$\rangle$P2]}}
                \texttt{See on minu KIRJELDUSKIRJELDUS: Ruudustik on 5x5. Esimeses reas on keskel X ja ääres tühjad ruudud. Teises reas on X esimeses, kolmandas ja viies veerus. Kolmandas, neljandas ja viiendas reas on X esimeses, teises, neljandas ja viiendas veerus, keskel on tühik.} \\
\texttt{Nüüd esitage küsimus, et teada saada minu ruudustiku uusi aspekte, mis võivad erineda teie ruudustikust. Alustage "KÜSIMUS:" ja ärge lisage midagi muud.} \\
            }
        }
    }
    & \\ \\

    \theutterance \stepcounter{utterance}  
    & & & \multicolumn{4}{p{0.6\linewidth}}{
        \cellcolor[rgb]{0.9,0.9,0.9}{
            \makecell[{{p{\linewidth}}}]{
                \texttt{\tiny{[GM$\langle$P2]}}
                \texttt{KÜSIMUS: Kas teie ruudustikus on X{-}tähed paigutatud sümmeetriliselt diagonaali või telje suhtes, või on muster asümmeetriline?} \\
            }
        }
    }
     \\ \\

    \theutterance \stepcounter{utterance}  
    & & & \multicolumn{2}{p{0.3\linewidth}}{
        \cellcolor[rgb]{0.9,0.9,0.9}{
            \makecell[{{p{\linewidth}}}]{
                \texttt{\tiny{[GM$|$GM]}}
                \texttt{continue} \\
            }
        }
    }
    & & \\ \\

    \theutterance \stepcounter{utterance}  
    & & \multicolumn{4}{p{0.6\linewidth}}{
        \cellcolor[rgb]{0.9,0.9,0.9}{
            \makecell[{{p{\linewidth}}}]{
                \texttt{\tiny{[P1$\langle$GM]}}
                \texttt{See on minu KIRJELDUSKIRJELDUS: Ruudustik on 5x5. X{-}tähed moodustavad kolmnurga, mis algab paremast ülemisest nurgast ja laieneb allapoole vasakule. Täpsemalt:} \\
\texttt{{-} 1. rida: 4 tühja ruutu, siis 1 X (positsioonil 1,5)} \\
\texttt{{-} 2. rida: 3 tühja ruutu, siis 2 X{-}d (positsioonidel 2,4 ja 2,5)} \\
\texttt{{-} 3. rida: 2 tühja ruutu, siis 3 X{-}d (positsioonidel 3,3, 3,4 ja 3,5)} \\
\texttt{{-} 4. rida: 1 tühi ruut, siis 4 X{-}d (positsioonidel 4,2 kuni 4,5)} \\
\texttt{{-} 5. rida: 5 X{-}d (kogu rida on X{-}tähtedega täidetud)} \\
\texttt{KÜSIMUS: Kas teie ruudustikus on X{-}tähed paigutatud sümmeetriliselt diagonaali või telje suhtes, või on muster asümmeetriline?Alustage oma vastust "VASTUS:" ja ärge lisage midagi muud.} \\
            }
        }
    }
    & \\ \\

    \theutterance \stepcounter{utterance}  
    & \multicolumn{4}{p{0.6\linewidth}}{
        \cellcolor[rgb]{0.9,0.9,0.9}{
            \makecell[{{p{\linewidth}}}]{
                \texttt{\tiny{[P1$\rangle$GM]}}
                \texttt{VASTUS: Asümmeetriline.} \\
            }
        }
    }
    & & \\ \\

    \theutterance \stepcounter{utterance}  
    & & & \multicolumn{2}{p{0.3\linewidth}}{
        \cellcolor[rgb]{0.9,0.9,0.9}{
            \makecell[{{p{\linewidth}}}]{
                \texttt{\tiny{[GM$|$GM]}}
                \texttt{continue} \\
            }
        }
    }
    & & \\ \\

    \theutterance \stepcounter{utterance}  
    & & \multicolumn{4}{p{0.6\linewidth}}{
        \cellcolor[rgb]{0.9,0.9,0.9}{
            \makecell[{{p{\linewidth}}}]{
                \texttt{\tiny{[P1$\langle$GM]}}
                \texttt{Nüüd esitage küsimus, et teada saada minu ruudustiku uusi aspekte, mis võivad erineda teie ruudustikust. Alustage "KÜSIMUS:" ja ärge lisage midagi muud.} \\
            }
        }
    }
    & \\ \\

    \theutterance \stepcounter{utterance}  
    & \multicolumn{4}{p{0.6\linewidth}}{
        \cellcolor[rgb]{0.9,0.9,0.9}{
            \makecell[{{p{\linewidth}}}]{
                \texttt{\tiny{[P1$\rangle$GM]}}
                \texttt{KÜSIMUS: Kas teie ruudustikus on X{-}tähed paigutatud nii, et need moodustavad ühtlase kolmnurga, mille tipp on paremal ülemisel nurgal?} \\
            }
        }
    }
    & & \\ \\

    \theutterance \stepcounter{utterance}  
    & & & \multicolumn{2}{p{0.3\linewidth}}{
        \cellcolor[rgb]{0.9,0.9,0.9}{
            \makecell[{{p{\linewidth}}}]{
                \texttt{\tiny{[GM$|$GM]}}
                \texttt{continue} \\
            }
        }
    }
    & & \\ \\

    \theutterance \stepcounter{utterance}  
    & & \multicolumn{4}{p{0.6\linewidth}}{
        \cellcolor[rgb]{0.9,0.9,0.9}{
            \makecell[{{p{\linewidth}}}]{
                \texttt{\tiny{[GM$\rangle$P2]}}
                \texttt{VASTUS: Asümmeetriline.} \\
\texttt{KÜSIMUS: Kas teie ruudustikus on X{-}tähed paigutatud nii, et need moodustavad ühtlase kolmnurga, mille tipp on paremal ülemisel nurgal?Alustage oma vastust "VASTUS:" ja ärge lisage midagi muud.} \\
            }
        }
    }
    & \\ \\

    \theutterance \stepcounter{utterance}  
    & & & \multicolumn{4}{p{0.6\linewidth}}{
        \cellcolor[rgb]{0.9,0.9,0.9}{
            \makecell[{{p{\linewidth}}}]{
                \texttt{\tiny{[GM$\langle$P2]}}
                \texttt{VASTUS: Ei, minu ruudustikus moodustavad X{-}tähed kolmnurga, mille tipp on paremal ülemisel nurgal, kuid see kolmnurk on täidetud ainult servadega, mitte täis. Teie kirjeldus viitab sellele, et X{-}d on hajutatud erinevalt, eriti keskosas, kus teie ruudustikus on keskel tühik, minu omas aga X{-}d.} \\
            }
        }
    }
     \\ \\

    \theutterance \stepcounter{utterance}  
    & & & \multicolumn{2}{p{0.3\linewidth}}{
        \cellcolor[rgb]{0.9,0.9,0.9}{
            \makecell[{{p{\linewidth}}}]{
                \texttt{\tiny{[GM$|$GM]}}
                \texttt{continue} \\
            }
        }
    }
    & & \\ \\

    \theutterance \stepcounter{utterance}  
    & & \multicolumn{4}{p{0.6\linewidth}}{
        \cellcolor[rgb]{0.9,0.9,0.9}{
            \makecell[{{p{\linewidth}}}]{
                \texttt{\tiny{[GM$\rangle$P2]}}
                \texttt{Nüüd esitage küsimus, et teada saada minu ruudustiku uusi aspekte, mis võivad erineda teie ruudustikust. Alustage "KÜSIMUS:" ja ärge lisage midagi muud.} \\
            }
        }
    }
    & \\ \\

    \theutterance \stepcounter{utterance}  
    & & & \multicolumn{4}{p{0.6\linewidth}}{
        \cellcolor[rgb]{0.9,0.9,0.9}{
            \makecell[{{p{\linewidth}}}]{
                \texttt{\tiny{[GM$\langle$P2]}}
                \texttt{KÜSIMUS: Kas teie ruudustiku teises reas on X{-}d veergudes 1, 3 ja 5, nagu kirjeldasite, või on see kirjeldus ekslik ja X{-}d asuvad hoopis teistes veergudes?} \\
            }
        }
    }
     \\ \\

    \theutterance \stepcounter{utterance}  
    & & & \multicolumn{2}{p{0.3\linewidth}}{
        \cellcolor[rgb]{0.9,0.9,0.9}{
            \makecell[{{p{\linewidth}}}]{
                \texttt{\tiny{[GM$|$GM]}}
                \texttt{continue} \\
            }
        }
    }
    & & \\ \\

    \theutterance \stepcounter{utterance}  
    & & \multicolumn{4}{p{0.6\linewidth}}{
        \cellcolor[rgb]{0.9,0.9,0.9}{
            \makecell[{{p{\linewidth}}}]{
                \texttt{\tiny{[P1$\langle$GM]}}
                \texttt{VASTUS: Ei, minu ruudustikus moodustavad X{-}tähed kolmnurga, mille tipp on paremal ülemisel nurgal, kuid see kolmnurk on täidetud ainult servadega, mitte täis. Teie kirjeldus viitab sellele, et X{-}d on hajutatud erinevalt, eriti keskosas, kus teie ruudustikus on keskel tühik, minu omas aga X{-}d.} \\
\texttt{KÜSIMUS: Kas teie ruudustiku teises reas on X{-}d veergudes 1, 3 ja 5, nagu kirjeldasite, või on see kirjeldus ekslik ja X{-}d asuvad hoopis teistes veergudes?Alustage oma vastust "VASTUS:" ja ärge lisage midagi muud.} \\
            }
        }
    }
    & \\ \\

    \theutterance \stepcounter{utterance}  
    & \multicolumn{4}{p{0.6\linewidth}}{
        \cellcolor[rgb]{0.9,0.9,0.9}{
            \makecell[{{p{\linewidth}}}]{
                \texttt{\tiny{[P1$\rangle$GM]}}
                \texttt{VASTUS: Jah, minu ruudustiku teises reas on X{-}d veergudes 1, 3 ja 5.} \\
            }
        }
    }
    & & \\ \\

    \theutterance \stepcounter{utterance}  
    & & & \multicolumn{2}{p{0.3\linewidth}}{
        \cellcolor[rgb]{0.9,0.9,0.9}{
            \makecell[{{p{\linewidth}}}]{
                \texttt{\tiny{[GM$|$GM]}}
                \texttt{continue} \\
            }
        }
    }
    & & \\ \\

    \theutterance \stepcounter{utterance}  
    & & \multicolumn{4}{p{0.6\linewidth}}{
        \cellcolor[rgb]{0.9,0.9,0.9}{
            \makecell[{{p{\linewidth}}}]{
                \texttt{\tiny{[P1$\langle$GM]}}
                \texttt{Nüüd esitage küsimus, et teada saada minu ruudustiku uusi aspekte, mis võivad erineda teie ruudustikust. Alustage "KÜSIMUS:" ja ärge lisage midagi muud.} \\
            }
        }
    }
    & \\ \\

    \theutterance \stepcounter{utterance}  
    & \multicolumn{4}{p{0.6\linewidth}}{
        \cellcolor[rgb]{0.9,0.9,0.9}{
            \makecell[{{p{\linewidth}}}]{
                \texttt{\tiny{[P1$\rangle$GM]}}
                \texttt{KÜSIMUS: Kas teie ruudustiku kolmandas reas on X{-}d veergudes 3, 4 ja 5, nagu kirjeldasite, või on see kirjeldus ekslik?} \\
            }
        }
    }
    & & \\ \\

    \theutterance \stepcounter{utterance}  
    & & & \multicolumn{2}{p{0.3\linewidth}}{
        \cellcolor[rgb]{0.9,0.9,0.9}{
            \makecell[{{p{\linewidth}}}]{
                \texttt{\tiny{[GM$|$GM]}}
                \texttt{continue} \\
            }
        }
    }
    & & \\ \\

    \theutterance \stepcounter{utterance}  
    & & \multicolumn{4}{p{0.6\linewidth}}{
        \cellcolor[rgb]{0.9,0.9,0.9}{
            \makecell[{{p{\linewidth}}}]{
                \texttt{\tiny{[GM$\rangle$P2]}}
                \texttt{VASTUS: Jah, minu ruudustiku teises reas on X{-}d veergudes 1, 3 ja 5.} \\
\texttt{KÜSIMUS: Kas teie ruudustiku kolmandas reas on X{-}d veergudes 3, 4 ja 5, nagu kirjeldasite, või on see kirjeldus ekslik?Alustage oma vastust "VASTUS:" ja ärge lisage midagi muud.} \\
            }
        }
    }
    & \\ \\

    \theutterance \stepcounter{utterance}  
    & & & \multicolumn{4}{p{0.6\linewidth}}{
        \cellcolor[rgb]{0.9,0.9,0.9}{
            \makecell[{{p{\linewidth}}}]{
                \texttt{\tiny{[GM$\langle$P2]}}
                \texttt{VASTUS: Ei, minu ruudustiku kolmandas reas on X{-}d veergudes 3, 4 ja 5, kuid lisaks on ka veergudes 1 ja 2 X{-}d, mitte ainult 3, 4 ja 5. Teie kirjeldus viitab sellele, et kolmandas reas on X{-}d ainult veergudes 3, 4 ja 5, mis on minu ruudustikust erinev.} \\
            }
        }
    }
     \\ \\

    \theutterance \stepcounter{utterance}  
    & & & \multicolumn{2}{p{0.3\linewidth}}{
        \cellcolor[rgb]{0.9,0.9,0.9}{
            \makecell[{{p{\linewidth}}}]{
                \texttt{\tiny{[GM$|$GM]}}
                \texttt{continue} \\
            }
        }
    }
    & & \\ \\

    \theutterance \stepcounter{utterance}  
    & & \multicolumn{4}{p{0.6\linewidth}}{
        \cellcolor[rgb]{0.9,0.9,0.9}{
            \makecell[{{p{\linewidth}}}]{
                \texttt{\tiny{[GM$\rangle$P2]}}
                \texttt{Nüüd esitage küsimus, et teada saada minu ruudustiku uusi aspekte, mis võivad erineda teie ruudustikust. Alustage "KÜSIMUS:" ja ärge lisage midagi muud.} \\
            }
        }
    }
    & \\ \\

    \theutterance \stepcounter{utterance}  
    & & & \multicolumn{4}{p{0.6\linewidth}}{
        \cellcolor[rgb]{0.9,0.9,0.9}{
            \makecell[{{p{\linewidth}}}]{
                \texttt{\tiny{[GM$\langle$P2]}}
                \texttt{KÜSIMUS: Kas teie ruudustiku neljandas reas on X{-}d veergudes 1, 2, 4 ja 5, nagu kirjeldasite, või on see kirjeldus ekslik?} \\
            }
        }
    }
     \\ \\

    \theutterance \stepcounter{utterance}  
    & & & \multicolumn{2}{p{0.3\linewidth}}{
        \cellcolor[rgb]{0.9,0.9,0.9}{
            \makecell[{{p{\linewidth}}}]{
                \texttt{\tiny{[GM$|$GM]}}
                \texttt{continue} \\
            }
        }
    }
    & & \\ \\

    \theutterance \stepcounter{utterance}  
    & & \multicolumn{4}{p{0.6\linewidth}}{
        \cellcolor[rgb]{0.9,0.9,0.9}{
            \makecell[{{p{\linewidth}}}]{
                \texttt{\tiny{[P1$\langle$GM]}}
                \texttt{VASTUS: Ei, minu ruudustiku kolmandas reas on X{-}d veergudes 3, 4 ja 5, kuid lisaks on ka veergudes 1 ja 2 X{-}d, mitte ainult 3, 4 ja 5. Teie kirjeldus viitab sellele, et kolmandas reas on X{-}d ainult veergudes 3, 4 ja 5, mis on minu ruudustikust erinev.} \\
\texttt{KÜSIMUS: Kas teie ruudustiku neljandas reas on X{-}d veergudes 1, 2, 4 ja 5, nagu kirjeldasite, või on see kirjeldus ekslik?Alustage oma vastust "VASTUS:" ja ärge lisage midagi muud.} \\
            }
        }
    }
    & \\ \\

    \theutterance \stepcounter{utterance}  
    & \multicolumn{4}{p{0.6\linewidth}}{
        \cellcolor[rgb]{0.9,0.9,0.9}{
            \makecell[{{p{\linewidth}}}]{
                \texttt{\tiny{[P1$\rangle$GM]}}
                \texttt{VASTUS: Jah, minu ruudustiku neljandas reas on X{-}d veergudes 1, 2, 4 ja 5.} \\
            }
        }
    }
    & & \\ \\

    \theutterance \stepcounter{utterance}  
    & & & \multicolumn{2}{p{0.3\linewidth}}{
        \cellcolor[rgb]{0.9,0.9,0.9}{
            \makecell[{{p{\linewidth}}}]{
                \texttt{\tiny{[GM$|$GM]}}
                \texttt{continue} \\
            }
        }
    }
    & & \\ \\

    \theutterance \stepcounter{utterance}  
    & & \multicolumn{4}{p{0.6\linewidth}}{
        \cellcolor[rgb]{0.9,0.9,0.9}{
            \makecell[{{p{\linewidth}}}]{
                \texttt{\tiny{[P1$\langle$GM]}}
                \texttt{Nüüd esitage küsimus, et teada saada minu ruudustiku uusi aspekte, mis võivad erineda teie ruudustikust. Alustage "KÜSIMUS:" ja ärge lisage midagi muud.} \\
            }
        }
    }
    & \\ \\

    \theutterance \stepcounter{utterance}  
    & \multicolumn{4}{p{0.6\linewidth}}{
        \cellcolor[rgb]{0.9,0.9,0.9}{
            \makecell[{{p{\linewidth}}}]{
                \texttt{\tiny{[P1$\rangle$GM]}}
                \texttt{KÜSIMUS: Kas teie ruudustiku viiendas reas on X{-}d veergudes 1, 2, 4 ja 5, nagu kirjeldasite, või on see kirjeldus ekslik?} \\
            }
        }
    }
    & & \\ \\

    \theutterance \stepcounter{utterance}  
    & & & \multicolumn{2}{p{0.3\linewidth}}{
        \cellcolor[rgb]{0.9,0.9,0.9}{
            \makecell[{{p{\linewidth}}}]{
                \texttt{\tiny{[GM$|$GM]}}
                \texttt{continue} \\
            }
        }
    }
    & & \\ \\

    \theutterance \stepcounter{utterance}  
    & & \multicolumn{4}{p{0.6\linewidth}}{
        \cellcolor[rgb]{0.9,0.9,0.9}{
            \makecell[{{p{\linewidth}}}]{
                \texttt{\tiny{[GM$\rangle$P2]}}
                \texttt{VASTUS: Jah, minu ruudustiku neljandas reas on X{-}d veergudes 1, 2, 4 ja 5.} \\
\texttt{KÜSIMUS: Kas teie ruudustiku viiendas reas on X{-}d veergudes 1, 2, 4 ja 5, nagu kirjeldasite, või on see kirjeldus ekslik?Alustage oma vastust "VASTUS:" ja ärge lisage midagi muud.} \\
            }
        }
    }
    & \\ \\

    \theutterance \stepcounter{utterance}  
    & & & \multicolumn{4}{p{0.6\linewidth}}{
        \cellcolor[rgb]{0.9,0.9,0.9}{
            \makecell[{{p{\linewidth}}}]{
                \texttt{\tiny{[GM$\langle$P2]}}
                \texttt{VASTUS: Ei, minu ruudustiku viiendas reas on X{-}d kõikides viies veerus (1, 2, 3, 4 ja 5), mitte ainult veergudes 1, 2, 4 ja 5.} \\
            }
        }
    }
     \\ \\

    \theutterance \stepcounter{utterance}  
    & & & \multicolumn{2}{p{0.3\linewidth}}{
        \cellcolor[rgb]{0.9,0.9,0.9}{
            \makecell[{{p{\linewidth}}}]{
                \texttt{\tiny{[GM$|$GM]}}
                \texttt{continue} \\
            }
        }
    }
    & & \\ \\

    \theutterance \stepcounter{utterance}  
    & & \multicolumn{4}{p{0.6\linewidth}}{
        \cellcolor[rgb]{0.9,0.9,0.9}{
            \makecell[{{p{\linewidth}}}]{
                \texttt{\tiny{[GM$\rangle$P2]}}
                \texttt{Nüüd tee otsus. Mida sa arvad: kas sinu ruudustik ja teine kirjeldatud ruudustik on samad?  Kirjuta "OTSUS: sama ruudustik", kui arvad, et need on sama ruudustik, või "OTSUS: erinevad ruudustikud", kui arvad, et need on erinevad ruudustikud. Ära lisa midagi muud.} \\
            }
        }
    }
    & \\ \\

    \theutterance \stepcounter{utterance}  
    & & & \multicolumn{4}{p{0.6\linewidth}}{
        \cellcolor[rgb]{0.9,0.9,0.9}{
            \makecell[{{p{\linewidth}}}]{
                \texttt{\tiny{[GM$\langle$P2]}}
                \texttt{OTSUS: erinevad ruudustikud} \\
            }
        }
    }
     \\ \\

    \theutterance \stepcounter{utterance}  
    & & & \multicolumn{2}{p{0.3\linewidth}}{
        \cellcolor[rgb]{0.9,0.9,0.9}{
            \makecell[{{p{\linewidth}}}]{
                \texttt{\tiny{[GM$|$GM]}}
                \texttt{continue} \\
            }
        }
    }
    & & \\ \\

    \theutterance \stepcounter{utterance}  
    & & & \multicolumn{2}{p{0.3\linewidth}}{
        \cellcolor[rgb]{0.9,0.9,0.9}{
            \makecell[{{p{\linewidth}}}]{
                \texttt{\tiny{[GM$|$GM]}}
                \texttt{success} \\
            }
        }
    }
    & & \\ \\

    \theutterance \stepcounter{utterance}  
    & & \multicolumn{4}{p{0.6\linewidth}}{
        \cellcolor[rgb]{0.9,0.9,0.9}{
            \makecell[{{p{\linewidth}}}]{
                \texttt{\tiny{[P1$\langle$GM]}}
                \texttt{VASTUS: Ei, minu ruudustiku viiendas reas on X{-}d kõikides viies veerus (1, 2, 3, 4 ja 5), mitte ainult veergudes 1, 2, 4 ja 5.} \\
\texttt{Nüüd tee otsus. Mida sa arvad: kas sinu ruudustik ja teine kirjeldatud ruudustik on samad?  Kirjuta "OTSUS: sama ruudustik", kui arvad, et need on sama ruudustik, või "OTSUS: erinevad ruudustikud", kui arvad, et need on erinevad ruudustikud. Ära lisa midagi muud.} \\
            }
        }
    }
    & \\ \\

    \theutterance \stepcounter{utterance}  
    & \multicolumn{4}{p{0.6\linewidth}}{
        \cellcolor[rgb]{0.9,0.9,0.9}{
            \makecell[{{p{\linewidth}}}]{
                \texttt{\tiny{[P1$\rangle$GM]}}
                \texttt{OTSUS: erinevad ruudustikud} \\
            }
        }
    }
    & & \\ \\

    \theutterance \stepcounter{utterance}  
    & & & \multicolumn{2}{p{0.3\linewidth}}{
        \cellcolor[rgb]{0.9,0.9,0.9}{
            \makecell[{{p{\linewidth}}}]{
                \texttt{\tiny{[GM$|$GM]}}
                \texttt{continue} \\
            }
        }
    }
    & & \\ \\

    \theutterance \stepcounter{utterance}  
    & & & \multicolumn{2}{p{0.3\linewidth}}{
        \cellcolor[rgb]{0.9,0.9,0.9}{
            \makecell[{{p{\linewidth}}}]{
                \texttt{\tiny{[GM$|$GM]}}
                \texttt{success} \\
            }
        }
    }
    & & \\ \\

\end{supertabular}
}

\end{document}
